\documentclass[
10pt, % Main document font size
a4paper, % Paper type, use 'letterpaper' for US Letter paper
oneside, % One page layout (no page indentation)
%twoside, % Two page layout (page indentation for binding and different headers)
headinclude,footinclude, % Extra spacing for the header and footer
BCOR5mm, % Binding correction
]{scrartcl}
\input{stru/structure.tex}
\hyphenation{Fortran hy-phen-ation} 

\title{\normalfont{FTML practical session 5}} 
\author{\spacedlowsmallcaps{}
} 
\date{\today}
\begin{document}
\renewcommand{\sectionmark}[1]{\markright{\spacedlowsmallcaps{#1}}}
\lehead{\mbox{\llap{\small\thepage\kern1em\color{halfgray} \vline}\color{halfgray}\hspace{0.5em}\rightmark\hfil}} 
\pagestyle{scrheadings}
\maketitle 
\setcounter{tocdepth}{2} 

\begin{figure}[htpb]
    \centering
    \includegraphics[width=1.1\linewidth]{knn_k=2_ntest=1000.pdf}
    \label{fig:risks}
\end{figure}

\tableofcontents

\section{\large\color{Blue}Curse of dimensionality and adaptivity to a low dimensional support}

\subsection{Nearest neighbors algorithms}

A nearest neighbors algorithm is a method of prediction, based on a dataset, but without any optimization. It consists in averaging the predictions of the nearest neighbors
 of a sample $x$ in a given dataset. See two examples below, where both input and output spaces are $\mathbb{R}$. The two example differ by the number of samples in the dataset.

 \begin{figure}[H]
 	\centering
 	\includegraphics[width=0.8\textwidth]{knn100}
 	\caption{knn}
 	\label{fig:knn}
 \end{figure}

 \begin{figure}[H]
 	\centering
 	\includegraphics[width=0.8\textwidth]{knn}
 	\caption{knn}
 	\label{fig:knn}
 \end{figure}

 A natural question would be: why do we not use this method all the time ? The answer is that as soon as the input space dimension is not very small (for instance not < 10), the number of samples that would be required to obtain a satisfactory error is simply too large.

\subsection{1D and 2D case}

\subsubsection{1D case}

We consider a $L$-Lipshitz-continuous function $f$, defined over $[0,1]$.

\begin{equation}
    \forall x,y\in [0,1], |f(x)-f(y)|\leq L|x-y|
\end{equation}

{\color{Purple}{Show that if we have $n$ samples, we can approximate $f$ with a nearest neighbor method with an error smaller than $\frac{L}{2n}$}}.

\subsubsection{2D case}

We consider a $L$-Lipshitz-continuous function $f$, now defined over $[0,1]^2$.

\begin{equation}
    \forall x,y\in [0,1]^2, |f(x)-f(y)|\leq L|x-y|
\end{equation}

{\color{Purple}{Show that if we have $n$ samples, we can approximate $f$ with a nearest neighbor method with an error smaller than $\frac{L}{\sqrt{2n}}$}}.

\subsection{Curse of dimensionality: general case}



In the general case, it is possible to show that if the target function $f^*$ is Lipshitz-continuous,
and if the $n$ samples in the dataset are on a lattice that divides $[0, 1]^d$ in
cubes of equal size, then the error made by a k-NN esimator can be upper bounded by a bound of the form $\mathcal{O}( n^{-1/d})$. 

Note that this is a rough bound. It is possible to have more subtle results,
and the constant in the $ \mathcal{O} $ depends on $d$ and on the Lipshitz
constant of the target function However, it is
not possible to have a bound that is "way better" (i.e. "way smaller") than
this bound. This means that if $d$ is large, then $1/d$ is small, and the decrease of the
error is in general very slow (meaning: in general, we should expect to have a decrease that is a slow as this one).

With additional hypotheses, it is also possible to show that in order to guarantee an error as small as some value $\epsilon>0$, 
the number of sampled needed  $n_{\epsilon}$ verifies

\begin{equation}
\label{eq:error}
    n_{\epsilon} \geq \frac{\epsilon^{-d}d^{d/2}}{\alpha^{d/2}} 
\end{equation}

where $\alpha > 0$ is a constant. $n_{\epsilon}$ hence grows exponentially with
$1/\epsilon$, and more than exponentially with $d$ which is a example of curse of dimensionality.
\\

We note that our target function is indeed Lipshitz continuous and thus respect the hypothesis to obtain equation $\ref{eq:error}$. If a function is not Lipshitz continuous, the situation is even worse !

\subsubsection{Simulation: uniformly distributed data in the input space}

We will study the influence of the input space dimension on the possibility to
approximate a function with a nearest neighbors approach in a specific setting:

\begin{itemize}
	\item input space: $ \mathcal{X} = [0, 1]^d$. Hence, here the
        dimensionality of the problem is $d$. The data will be uniformly
        sampled in the input space.
	\item output space: $\mathcal{Y} = \mathbb{R}$
	\item toy function to approximate: $f: x\in \mathbb{R}^d\mapsto ||x||$ (euclidean norm).
\end{itemize}

{\color{Purple}{Run a simulation that computes the test error of a kNN
estimator, for different values of $n$ and $d$, in order to observe the curse
of dimensionality. The dataset should be uniformly sampled in the input space.
You can experiment with the number of neighbors used of $k$, for instance $2$}}.
\\

You can use the template file \textbf{knn\textunderscore uniform.py}. Note that as we know the function that we try to approximate, it is possible to perfeclty evaluate the error made by the k-NN algorithm for each sample.
For instance, you can observe images like the one in figure 
\ref{fig:glu}. In the simulation that generated the figure, the samples (both
the dataset and the test samples) were
drawn randomly in $ \mathcal{X}$, with a uniform distribution (hence the dataset is
not on a regular lattice, although you can do use a regular lattice in your
simulation).

\begin{figure}[H]
    \centering
    \includegraphics[width=1.1\linewidth]{knn_k=2_ntest=1000.pdf}
    \caption{Curse of dimensionality. We observe that the logarithm of the
    error is linear in the logarithm of the number of samples, with a slope
    that depends on the dimension $d$. When $d$ is large, the problem is that
    the curve is very flat, which means that the error decreases very slowly.}
    \label{fig:glu}
\end{figure}

%\subsubsection{\large\color{Periwinkle}Template}
%
%You can use \textbf{{knn\textunderscore 1.py}}, and edit the function
%\textbf{{predict()}}, that implements the k-NN algorithm.

\subsection{Adaptivity}

However, the situation is not always that catastrophic in high dimensions and we will illustrate this with an example. We keep exactly the same setting, but now the input dataset $X$ is given and fixed, and we use $d=30$.
\\


{\color{Purple}{
Perform a k-NN estimation using the data stored in
\textbf{{data\textunderscore knn/}} and use a varying number of samples $n$
that you use from the dataset. Plot the error as a function of $n$.}}

You can use the template file \textbf{knn\textunderscore adaptivity.py}. You should observe something like figure \ref{fig:2}. We note that the learning rate is way better than $ \mathcal{O} (n^{-1/d})$.

{\color{Purple}{
\begin{itemize}
    \item What could be an explanation of this phenomenon ?
    \item How could we test this hypothesis ?
\end{itemize}
}}

\begin{figure}[htpb]
    \centering
    \includegraphics[width=0.8\linewidth]{x_test_knn_k=2}
    \caption{Nearest neighbors estimation on a different dataset.}
    \label{fig:2}
\end{figure}

%\subsubsection{\large\color{Periwinkle}Template files}
%\begin{itemize}
%    \item You can start from \textbf{{knn\textunderscore 2.py}} and import the
%        relevant functions from the previous file.
%    \item \textbf{{x\textunderscore data\textunderscore .npy}} and  \textbf{{y\textunderscore data\textunderscore .npy}}: dataset that
%        is used by the k-NN algorithm.
%    \item \textbf{{x\textunderscore data\textunderscore test.npy}} : data on
%        which we estimate using the k-NN algorithm.
%\end{itemize}

\end{document}
